% SHARD-ID: AQM-55-DERHAM-CSS-MORE-F3-EXAMPLES
% SHARD-TITLE: More \(\mathbb F_3\) Examples for the de Rham CSS Prototype
% SHARD-SUMMARY: Computes the finite de Rham CSS prototype for five finite \(\mathbb F_3\)-coordinate rings cut out by linear and quadratic equations.
% SHARD-SUMMARY: Separates the squarefree reduced cases, where \(\Omega^1\) vanishes, from nilpotent quadratic cases with derivative stabilizers.
% SHARD-SUMMARY: Lists the auxiliary \(C^1\)-qutrits and explicit \(X\)- and \(Z\)-stabilizer generators inside this prototype.
% SHARD-KEYWORDS: F3, de Rham, CSS, stabilizer, Kahler differentials, finite rings, quadratic, nilpotent
\section{More \texorpdfstring{\(\mathbb F_3\)}{F3} Examples for the de Rham CSS Prototype}
\label{sec:derham-css-more-f3-examples}
\subsection{Status, Sources, and Common Rule}
This shard is a follow-up calculation for the finite chain-complex prototype in
AQM-54 and convention \texttt{CONVENTIONS.md} (ao), not the canonical
tangent-cotangent Weyl quantisation.  The prototype chooses a finite
\(C^1=\Omega^1_{A/k}\) basis, attaches auxiliary qudits to it, and tests whether
the de Rham differential selects commuting Pauli operators.  All rings below are
finite-dimensional coordinate rings over
\[
  k=\mathbb F_3.
\]
The derivative input is the module of K\"ahler differentials
\(\Omega^1_{A/k}\); the prototype uses auxiliary \(C^1\)-qutrits and selects CSS
support spaces
\[
  R_X=\operatorname{im}(d_0:C^0\to C^1),\qquad
  R_Z=\operatorname{im}(d_1^T:C^2{}^*\to C^1{}^*),
\]
with \(d_1d_0=0\) giving the CSS commutation check.  The source anchors are the
same as AQM-54: Stacks algebra lines 33793--33872, 34103--34120, 34310--34340,
and 34388--34524.  AQM-08 supplies the CSS stabilizer dictionary.
The common one-variable calculation is this.  Let
\(A=k[t]/(f)\), \(\tau=t\bmod(f)\).  For \(S=k[t]\), the polynomial
differential rule gives \(\Omega^1_{S/k}=S\,dt\).  The quotient exact
sequence sends \(f\) to \(df=f'(t)\,dt\), hence
\[
  \Omega^1_{A/k}\cong A\,d\tau/(f'(\tau)d\tau).
\]
Because \(\Omega^1_{A/k}\) is generated by one symbol in all one-variable
examples, \(\Omega^2_{A/k}=0\) there and \(R_Z=0\).
\subsection{Lamport Registers for the Five Examples}
\[
\begin{array}{rcl}
E0&:&A_0=k[t]/(t)\cong k.\\
E1&:&A_1=k[t]/(t^2-1),\quad\text{the split quadratic two-point ring.}\\
E2&:&A_2=k[t]/(t^2+1),\quad\text{the irreducible quadratic field point.}\\
E3&:&A_3=k[t]/(t^2),\quad\text{the double point at }0.\\
E4&:&A_4=k[u,v]/(u^2,v^2),\quad\text{the quadratic fat plane point.}
\end{array}
\]
The examples are ordered so that \(E0,E1,E2\) test the reduced squarefree
cases, while \(E3,E4\) test nilpotent quadratic directions.
\subsection{E0: The Linear Point}
Let \(A_0=k[t]/(t)\), and let \(\tau=0\) be the image of \(t\).  Here
\(f=t\), \(f'=1\).  The common rule gives
\[
  \Omega^1_{A_0/k}=A_0\,d\tau/(1\cdot d\tau)=0.
\]
Thus
\[
  C^1=0,\qquad R_X=0,\qquad R_Z=0.
\]
The de Rham CSS test selects no qutrit and no stabilizer generator.  This is
the same result as AQM-54 Example 1, now recorded as the linear hypersurface
case \(t=0\).
\subsection{E1: The Split Quadratic Two-Point Ring}
Let
\[
  A_1=k[t]/(t^2-1),\qquad \tau=t\bmod(t^2-1).
\]
The relation is \(\tau^2=1\).  The defining polynomial and derivative are
\[
  f=t^2-1,\qquad f'=2t.
\]
In \(A_1\),
\[
  (2\tau)(2\tau)=4\tau^2=\tau^2=1,
\]
because \(4=1\) in \(k\).  Hence \(2\tau\) is a unit.  The quotient relation
is therefore
\[
  2\tau\,d\tau=0,
\]
and multiplication by the inverse of \(2\tau\) gives
\[
  d\tau=0.
\]
Consequently
\[
  \Omega^1_{A_1/k}=0.
\]
The de Rham CSS data are
\[
  C^1=0,\qquad R_X=0,\qquad R_Z=0.
\]
Thus the split reduced quadratic coordinate ring contributes two reduced
points but no intrinsic derivative qutrit.  In this test, reduced squarefree
finite geometry has kinematics from the Weyl layer but no derivative
stabilizer selected by \(\Omega^1\).
\subsection{E2: The Irreducible Quadratic Field Point}
Let
\[
  A_2=k[t]/(t^2+1),\qquad \tau=t\bmod(t^2+1).
\]
Since \(0^2+1=1\), \(1^2+1=2\), and \(2^2+1=5=2\) in \(k\), the polynomial
\(t^2+1\) has no \(k\)-root.  The ring \(A_2\) is the quadratic residue-field
extension represented by the closed point \(t^2+1\).  Its relation is
\[
  \tau^2=-1=2.
\]
Again
\[
  f=t^2+1,\qquad f'=2t.
\]
In \(A_2\),
\[
  (2\tau)\tau=2\tau^2=2\cdot2=4=1.
\]
Thus \(2\tau\) is a unit, now with inverse \(\tau\).  The quotient relation
\[
  2\tau\,d\tau=0
\]
forces
\[
  d\tau=0,\qquad \Omega^1_{A_2/k}=0.
\]
Therefore
\[
  C^1=0,\qquad R_X=0,\qquad R_Z=0.
\]
The irreducible quadratic closed point behaves like the split squarefree
quadratic for this derivative test: separability kills first differentials.
\subsection{E3: The Quadratic Double Point}
Let
\[
  A_3=k[t]/(t^2),\qquad \epsilon=t\bmod(t^2).
\]
Then \(A_3=k\oplus k\epsilon\) and \(\epsilon^2=0\).  Here
\[
  f=t^2,\qquad f'=2t.
\]
The common rule gives
\[
  \Omega^1_{A_3/k}=A_3\,d\epsilon/(2\epsilon\,d\epsilon).
\]
Since \(2\) is a unit in \(k\), this is
\[
  \Omega^1_{A_3/k}=A_3\,d\epsilon/(\epsilon\,d\epsilon)
  \cong k\,d\epsilon.
\]
The universal derivation is computed on the basis \(1,\epsilon\):
\[
  d_0(a+b\epsilon)=b\,d\epsilon.
\]
Thus
\[
  C^1=k\,d\epsilon,\qquad
  R_X=\operatorname{im}d_0=k\,d\epsilon.
\]
Because \(C^1\) has one generator, \(C^2=\Omega^2_{A_3/k}=0\), so
\[
  R_Z=0.
\]
In the prototype, with one auxiliary qutrit labelled by \(e_1=d\epsilon\), the
selected stabilizer is
\[
  X(e_1).
\]
There is no \(Z\)-check.  The stabilizer rank is \(1\) on one qutrit, so the
all-\(+1\) stabilizer space has dimension \(3^{1-1}=1\).
\subsection{E4: The Quadratic Fat Plane Point}
Let
\[
  A_4=k[u,v]/(u^2,v^2),
\]
where \(u,v\) also denote the residue classes.  A \(k\)-basis is
\[
  1,\quad u,\quad v,\quad uv,
\]
with \(u^2=v^2=0\).  Start from \(S=k[U,V]\).  The polynomial differential
rule gives
\[
  \Omega^1_{S/k}=S\,dU\oplus S\,dV.
\]
The quotient ideal is \((U^2,V^2)\).  The two differential relations are
\[
  d(U^2)=2U\,dU,\qquad d(V^2)=2V\,dV.
\]
Passing to \(A_4\), and using that \(2\) is a unit, gives
\[
  u\,du=0,\qquad v\,dv=0.
\]
Therefore
\[
  \Omega^1_{A_4/k}
  =
  (A_4\,du\oplus A_4\,dv)/(u\,du,\ v\,dv).
\]
The \(A_4\,du\) part has basis \(du,v\,du\), since \(u\,du=0\) and hence
\(uv\,du=v(u\,du)=0\).  The \(A_4\,dv\) part has basis \(dv,u\,dv\), since
\(v\,dv=0\) and hence \(uv\,dv=u(v\,dv)=0\).  Set
\[
  e_1=du,\quad e_2=v\,du,\quad e_3=dv,\quad e_4=u\,dv.
\]
Thus \(C^1=\Omega^1_{A_4/k}\cong k^4\), so the prototype has four auxiliary
qutrits.
The derivation \(d_0:A_4\to C^1\) is fixed on the basis:
\[
\begin{array}{rcl}
d_0(1)&=&0,\\
d_0(u)&=&e_1,\\
d_0(v)&=&e_3,\\
d_0(uv)&=&u\,dv+v\,du=e_4+e_2.
\end{array}
\]
Hence
\[
  R_X=\langle e_1,\ e_3,\ e_2+e_4\rangle_k.
\]
Now compute \(C^2=\Omega^2_{A_4/k}\).  It is generated by
\[
  \omega=du\wedge dv.
\]
Wedging \(u\,du=0\) with \(dv\) gives \(u\omega=0\), and wedging
\(v\,dv=0\) with \(du\) gives \(v\omega=0\).  Hence
\[
  C^2=\Omega^2_{A_4/k}\cong k\,\omega.
\]
The de Rham differential \(d_1:C^1\to C^2\) is
\[
\begin{array}{rcl}
d_1(e_1)&=&d(du)=0,\\
d_1(e_2)&=&d(v\,du)=dv\wedge du=-\omega=2\omega,\\
d_1(e_3)&=&d(dv)=0,\\
d_1(e_4)&=&d(u\,dv)=du\wedge dv=\omega.
\end{array}
\]
With the ordered bases \(e_1,e_2,e_3,e_4\) and \(\omega\), the matrix of
\(d_1\) is
\[
  \begin{bmatrix}0&2&0&1\end{bmatrix}.
\]
Thus the \(Z\)-support selected by the transpose is
\[
  R_Z=\langle 2e_2+e_4\rangle_k.
\]
The CSS orthogonality is visible without abbreviation:
\[
  e_1\cdot(2e_2+e_4)=0,\quad
  e_3\cdot(2e_2+e_4)=0,\quad
  (e_2+e_4)\cdot(2e_2+e_4)=2+1=0.
\]
Therefore the selected stabilizer generators on the four qutrits are
\[
  X_1,\qquad X_3,\qquad X_2X_4,\qquad Z_2^2Z_4.
\]
Their rank is \(4\), so the all-\(+1\) stabilizer space has dimension
\[
  3^{4-4}=1.
\]
\subsection{End Summary: Prototype Qudits and Stabilizers}
For \(C^1\)-basis vector \(e_i\), write \(X_i=X(e_i)\) and
\(Z_i=Z(e_i)\).  The de Rham CSS test gives:
{\scriptsize
\[
\begin{array}{c|c|c|c}
\text{example} & C^1\text{-qutrits} & X\text{-checks from }R_X
  & Z\text{-checks from }R_Z\\ \hline
E0:\ k[t]/(t) & \text{none} & \text{none} & \text{none}\\
E1:\ k[t]/(t^2-1) & \text{none} & \text{none} & \text{none}\\
E2:\ k[t]/(t^2+1) & \text{none} & \text{none} & \text{none}\\
E3:\ k[t]/(t^2) & e_1=d\epsilon & X_1 & \text{none}\\
E4:\ k[u,v]/(u^2,v^2) &
  e_1=du,\ e_2=vdu,\ e_3=dv,\ e_4=udv &
  X_1,\ X_3,\ X_2X_4 &
  Z_2^2Z_4
\end{array}
\]
}
Thus this prototype detects nilpotent or singular first-order algebra in
\(C^1\) and is silent on these squarefree finite reduced examples.  The
canonical tangent-Weyl reading must instead pass through the pointwise spaces
\(T_PX\oplus T_P^*X\) and their Weyl representations.
